\section{Conclusion}
\label{sec:conclusion}


Current embodied intelligence systems face significant challenges in human-robot interaction, including high latency, rigid interaction patterns, and an inability to handle interruptions gracefully. The VITA-VLA framework addresses these limitations through its innovative dual-model architecture and special token-based control flow, taking an important step toward achieving nearly real-time, interruptible, and concurrent human-robot interaction.

Our experimental results demonstrate that the framework significantly reduces interaction latency and interruption response time while enabling natural concurrent behaviors. The "model-as-controller" paradigm, where large language models generate their own control tokens, represents a novel approach to tightly coupling AI reasoning with system behavior.

We envision that this work will inspire future research toward more natural and intelligent human-robot collaboration, ultimately realizing the goal of seamless embodied AI assistants that can adapt and respond to human needs in real-time dynamic environments.